\section{Safety: what in the bill actually protects patients?}
\label{sec:safety}

The third question is the one a member will be asked on local television: what, in
plain terms, keeps a robot from hurting someone? The answer must be concrete, and it
must be in the statute rather than in a press release.

The mechanism is verification before generation. Before any robot-patient
interaction code is generated or executed, the system's proposed action is examined
by a ten-gate VVUQ suite, and the action is resolved to ACCEPT and proceed under
audit, ESCALATE to a qualified human, or BLOCK before it can reach the patient. The
decision rule in section 515D is strict: a single failing dimension forces a BLOCK,
any ambiguity or divergence forces an ESCALATE, and the verification fraction, the
share of interactions that must clear every gate, is fixed at 1.0
\cite{Kawchak2026HR9510Bill}. \autoref{fig:vbg} shows the pipeline.

\begin{psgfig}
\node[psgone] (p) at (0,0) {Propose\\action};
\node[psgtwo] (r) at (3.5,0) {Independent\\review};
\node[psgdec] (g) at (7.0,0) {Ten\\gates};
\node[psgevid] (a) at (10.6,1.2) {ACCEPT\\under audit};
\node[psgthree] (e) at (10.6,0) {ESCALATE\\to a human};
\node[psgcaut] (b) at (10.6,-1.2) {BLOCK\\before patient};
\draw[psgedge] (p)--(r); \draw[psgedge] (r)--(g);
\draw[psgedge] (g)--(a.south west); \draw[psgedge] (g)--(e.west);
\draw[psgedge] (g)--(b.north west);
\draw[psgedged] (e) to[out=180,in=-10,looseness=0.5] (g.south east);
\end{psgfig}
\psgcaption{\textbf{Figure 4.} Verification before generation: every proposed action
is gated to ACCEPT, ESCALATE, or BLOCK before the patient; the colored TikZ
reproduction of catalog diagram 03.}

What makes the floor credible is that the bill does not invent its own rigor. Each
of the ten gates is bound to one or more published consensus standards, with a
minimum validation agreement and a maximum coefficient of variation, and three of the
gates, the vascular no-fly volume, the shared-room human collaboration, and the
fault and emergency-stop behavior, are hard catastrophe predicates that must hold at
an agreement of 1.00 without exception \cite{asme2018vv40, iso14971, isots15066,
ul4600, fda2021samd}. \autoref{tab:gates} reproduces the schedule. That the suite is
not theoretical is shown by the author's assurance lineage, in which a surgical
humanoid was cleared only after the same gate discipline was applied
\cite{Kawchak2026MobilePancreaticCancerUnitree, Kawchak2026VVUQOncologyClinicalTrial}.

\needspace{16\baselineskip}\begin{table}[H]
\footnotesize
\begin{tabularx}{\textwidth}{@{}L{0.6cm} L{3.5cm} C{1.5cm} C{1.0cm} Y@{}}
\toprule
\textbf{\#} & \textbf{Gate} & \textbf{Agree.} & \textbf{CV} & \textbf{Bound standards}\\
\midrule
1 & Bimanual hand-eye servo & 0.97 & 0.08 & IEC 80601-2-77; ISO 9283; ASME V\&V 40\\
2 & Dexterous finger force & 0.95 & 0.10 & IEC 80601-2-77; ISO/TS 15066\\
3 & Whole-body balance & 0.98 & 0.06 & ISO 13482; IEC 60601-1\\
4 & Autonomous plan correctness & 0.95 & 0.10 & UL 4600; IEEE 7009; IEC 62304\\
5 & Instructed grasp and handover & 0.96 & 0.10 & IEC 80601-2-77; ISO 9283\\
6 & Vascular no-fly volume (hard) & 1.00 & 0.05 & IEC 80601-2-77; ISO 14971\\
7 & Bimanual suturing anastomosis & 0.96 & 0.08 & IEC 80601-2-77; Fistula Risk Score\\
8 & Perception scene understanding & 0.95 & 0.10 & ASME V\&V 40; IEC 62304\\
9 & Shared-room collaboration (hard) & 1.00 & 0.06 & ISO/TS 15066; ISO 10218-1; ISO 13482\\
10 & Fault and emergency stop (hard) & 1.00 & 0.05 & IEC 60601-1; ISO 13849-1; IEEE 7009\\
\bottomrule
\end{tabularx}
\caption{The ten-gate threshold schedule of section 515D(c): each gate's minimum
validation agreement, maximum coefficient of variation, and bound consensus
standards. The three hard catastrophe predicates must hold at agreement 1.00.}
\label{tab:gates}
\end{table}

Safety is answered by a mechanism, not a promise: a ten-gate floor bound to standards
the relevant professions already trust, with three predicates that can stop the
system cold, all required before the code is even generated.
